\section{Conclusão}
\label{sec:conclusao}

Este trabalho investigou como um sistema inteligente pode monitorar a postura
de quem trabalha sentado e devolver, em tempo real, um aviso que leve à
autocorreção. Para isso, desenvolveu-se um protótipo completo e autônomo:
uma unidade vestível com IMU, posicionada logo abaixo do pescoço, e uma
cadeira de escritório com oito sensores de força resistivos, ligadas por um
enlace sem fio de baixa latência. A cadeira combina as duas fontes de
informação, decide quando a postura é prejudicial e envia o alerta, que a
unidade vestível converte em vibração.

Os resultados permitem três conclusões principais. Primeiro, a referência de
postura ereta varia entre pessoas e entre colocações do sensor em magnitude
comparável ao próprio desvio a detectar; por isso, a calibração individual
por botão, que define o ereto de cada usuário como origem, é condição para o
funcionamento, e não um refinamento. Segundo, sensores de força de baixo
custo, inadequados como medidores de força, mostraram-se suficientes como
indicadores de contato e de pressão relativa, desde que comparados com a
linha de base do próprio usuário. Terceiro, a fusão dessas leituras com a
orientação do tronco resolveu as falhas mais frequentes da abordagem baseada
apenas no IMU: deixou de alertar quando o usuário reclina com o quadril
apoiado ou se levanta da cadeira, e passou a detectar o escorregamento do
quadril mesmo com pouca inclinação do tronco. Nas capturas com oito
participantes, avaliadas deixando cada um de fora do ajuste, o sistema alertou
em 58 de 62 segmentos de postura prejudicial, com alerta em $3\,\%$ do tempo
em posturas aceitáveis; na demonstração pública, a postura prejudicial foi
detectada em praticamente todos os casos observados.

Esses resultados indicam a viabilidade técnica de detectar posturas
prejudiciais com a combinação proposta, mas não demonstram efeito sobre a
saúde. O conjunto de dados é pequeno, as posturas foram instruídas, e
permanecem limitações conhecidas: o movimento isolado da cabeça, a fixação do
colar ao longo da sessão, a interação entre vibração e sensor e o consumo de
energia.

\subsection{Trabalhos futuros}
\label{sec:trabalhos-futuros}

\begin{itemize}[nosep]
  \item \textbf{Validação ampliada:} novas sessões com mais participantes,
        posturas naturais e sutis, ordem embaralhada e uso prolongado em
        trabalho real, com o motor ativo e a correção da interferência do
        atuador já embarcada.
  \item \textbf{Padrão de vibração:} substituir os pulsos contínuos por um
        padrão com pausas crescentes, reduzindo incômodo e habituação, e
        avaliar a aceitação pelos usuários.
  \item \textbf{Energia:} medir corrente e autonomia de cada unidade; reduzir
        as retransmissões do enlace, desligar o rádio entre envios e baixar a
        taxa para $5$--$10\,Hz$, compatível com a dinâmica da postura.
  \item \textbf{Robustez da referência:} compensar a deriva lenta do colar
        quando o padrão dos FSR coincide com a linha de base, e aprimorar a
        fixação do invólucro.
  \item \textbf{Cabeça e pescoço:} investigar outra posição para o sensor ou
        um segundo IMU para detectar a projeção da cabeça à frente.
  \item \textbf{Assento:} avaliar mais pontos de pressão no assento para
        tratar pernas cruzadas e reduzir a dependência da espuma e do
        posicionamento de cada cadeira.
  \item \textbf{Efeito sobre o usuário:} estudo longitudinal que meça mudança
        de hábito postural e sintomas osteomusculares, por exemplo com o
        questionário nórdico, antes e depois do uso.
\end{itemize}
